\section{Tools Menu}
\label{sec:tools}

The \menu{Tools} menu collects supporting calculators that are useful
during CLS planning and analysis but do not belong to the main
four-tab workflow. Each tool opens in its own window and stays open
across tab switches, so values can be cross-checked while editing the
\app{Estimate} or \app{Pre-Analysis} forms.

\subsection{Schmidt Moment Calculator}

\menu{Tools $\rightarrow$ Schmidt Moment Calculator} computes
single-particle and two-particle Schmidt magnetic moments
(Figure~\ref{fig:tools_schmidt}). Choose \app{Single particle} for
odd-mass nuclei and supply the unpaired nucleon's type
(\code{p}/\code{n}) and orbital
(\code{1s1/2}, \code{1p3/2}, \dots, \code{1f7/2}, etc.); choose
\app{Two particles} for odd-odd nuclei and supply two orbital
configurations together with the resulting coupled spin $I$. The
\app{$g_s$ quenching factor} (default 1.0) lets the user apply a
literature quenching to the bare-nucleon spin $g$-factor. The result
is the predicted $\mu$ in nuclear magnetons, the corresponding
$g$-factor, and a one-line description of the configuration that is
also written into the Estimate log when Schmidt moments are used
there (Section~\ref{sec:estimate}).

\begin{figure}[hbtp]
    \centering
    \includegraphics[width=0.55\textwidth]{fig_tools_schmidt.png}
    \caption{Schmidt Moment Calculator with a proton in the
        $1s_{1/2}$ orbital --- the bare-nucleon prediction for a
        single odd proton.}
    \label{fig:tools_schmidt}
\end{figure}

\subsection{Shell Configuration Plotter}

\menu{Tools $\rightarrow$ Shell Configuration Plotter} draws the
nucleon shell-filling diagram for any nucleus
(Figure~\ref{fig:tools_shell}). Enter $Z$ and $N$ in the top-left
panel; the element symbol and mass number are derived automatically.
The plot shows protons (red, left) and neutrons (blue, right)
placed into spherical shell-model ($jj$-coupling) orbitals --- not a
deformed Nilsson diagram --- with magic-number gaps marked between
shells. The \app{Shell Model Results} block below the inputs
prints the resulting proton and neutron configuration strings and
the predicted ground-state $J^{\pi}$ from the Nordheim coupling
rule. \app{Plot Options} controls the depth of the shell ladder
shown (\app{From shell} / \app{To shell}), labels, vacancy markers,
$n$-quantum numbers, and circle size.

\begin{figure}[hbtp]
    \centering
    \includegraphics[width=0.95\textwidth]{fig_tools_shell.png}
    \caption{Shell Configuration Plotter for \textsuperscript{72}Ge
        ($Z=32$, $N=40$). Both proton and neutron shells are filled
        through the $1f_{5/2}$ orbital; the predicted $J^{\pi}$ is
        $0^+$ as expected for an even-even nucleus.}
    \label{fig:tools_shell}
\end{figure}

\subsection{Unit Converter}

\menu{Tools $\rightarrow$ Unit Converter} converts a
laser-spectroscopy quantity into every other unit it is likely to be
expressed in (Figure~\ref{fig:tools_unit_converter}). Enter a value
and pick the input unit (nm, cm$^{-1}$, MHz, GHz, THz, eV); the
converted values appear immediately in all six units. Useful when
toggling between the laser-control software (often nm), satlas2
parameters (MHz), and published transition tables (cm$^{-1}$).

\begin{figure}[hbtp]
    \centering
    \includegraphics[width=0.45\textwidth]{fig_tools_unit_converter.png}
    \caption{Unit Converter showing the second-harmonic of an
        Nd:YAG laser ($532$~nm) expressed in every supported unit.}
    \label{fig:tools_unit_converter}
\end{figure}

\subsection{SHG Crystal Angle Calculator}

\menu{Tools $\rightarrow$ SHG Crystal Angle Calculator} computes the
Type-I (o+o $\rightarrow$ e) phase-matching angle for second-harmonic
generation in nonlinear crystals, from Sellmeier dispersion
(Figure~\ref{fig:tools_shg}). Pick the crystal (\app{BBO},
\app{KDP}, \app{LBO}) and the calculation direction ---
\app{Wavelength $\rightarrow$ Angle} for typical use, or
\app{Angle $\rightarrow$ Wavelength} for diagnostics on a fixed
crystal mount. The result reports the phase-matching angle (or the
matched fundamental) and the SHG wavelength; if the crystal cannot
phase-match the requested line the result reads \app{No phase
matching}. Only Type-I matching is implemented.

\begin{figure}[hbtp]
    \centering
    \includegraphics[width=0.5\textwidth]{fig_tools_shg.png}
    \caption{SHG Crystal Angle Calculator: BBO Type-I phase-matching
        angle ($22.78^{\circ}$) for the second harmonic of 1064~nm.}
    \label{fig:tools_shg}
\end{figure}

\subsection{SFG / DFG Calculator}

\menu{Tools $\rightarrow$ SFG / DFG Calculator} combines two input
wavelengths into a sum-frequency (SFG) or difference-frequency (DFG)
output. Each input takes its own unit (nm, cm$^{-1}$, MHz, GHz, THz,
eV), the combination is done in wavenumber space (additive in photon
energy), and the generated wavelength is reported simultaneously in
all six units. SHG is simply the SFG special case with two equal
inputs.

\subsection{Quick Plot}

\menu{Tools $\rightarrow$ Quick Plot} is a lightweight data viewer
that lives outside the analysis pipeline
(Figure~\ref{fig:tools_quick_plot}). The \app{Load File} tab opens a
delimited file (auto-detected or explicit delimiter) and lists its
columns; \app{Enter Data} is a built-in spreadsheet that also accepts
pasted tab/comma/semicolon/whitespace-separated data. Pick \app{X}
and \app{Y} columns from the drop-downs, choose a chart type
(\app{Line}, \app{Scatter}, \app{Step}, or \app{Bar}), and click
\app{Plot}. Title, axis labels, grid and legend toggles, and a CSV /
TSV save round out the tool; the standard matplotlib toolbar provides
pan, zoom, and image export --- useful for sanity-checking
calibration files, ion-source trends, or any tabular data adjacent to
a CLS measurement.

\subsection{NIST ASD Browser}

\menu{Tools $\rightarrow$ NIST ASD Browser\dots} is a full in-app
browser for the NIST Atomic Spectra Database, plus a multi-step
laser excitation \emph{scheme finder}. Enter a spectrum in NIST
notation (e.g.\ \code{Si I}, \code{Ca II}) and click \app{Fetch}:
data is loaded from a local cache when present, otherwise downloaded
once from NIST and cached under \code{settings/nist\_cache}, so
subsequent sessions work offline. NIST data is fetched in vacuum;
the \app{Air}/\app{Vacuum} toggle converts locally with the standard
NIST dispersion formula.

\begin{description}
    \item[\app{Levels} tab.] Filterable level table (configuration,
        term, energy window, \app{Metastables only}) with three plot
        views: a Grotrian-style \app{Level diagram} (one column per
        configuration; bar opacity encodes the radiative lifetime
        $\tau = 1/\Sigma A_{ki}$, and hovering a level shows its
        energy, uncertainty, configuration/term/$J$, and $\tau$),
        \app{Decay channels} of the selected level (one arrow per
        channel with $\lambda$ and branching ratio --- pick your
        detection channel here), and \app{Lifetimes vs energy}
        (metastable candidates jump out). \app{Add as starting
        level} sends the selected level to the Scheme Finder.
    \item[\app{Lines} tab.] Filterable transition table
        ($\lambda$ range, minimum $A_{ki}$, allowed / forbidden,
        involves-level) with a \app{Stick spectrum}
        ($\lambda$ vs $\log_{10} A_{ki}$) and a \app{Line density}
        histogram (spectral congestion at a glance), plus manual
        axis-limit controls.
    \item[\app{Scheme Finder} tab.] Searches 1--3-step excitation
        schemes from chosen (or auto-discovered ground + metastable)
        starting levels. Lasers are either named systems with a
        centre $\pm$ range and a role (\app{PUMP} / \app{PROBE} /
        \app{ANY}; each used at most once per scheme) or one broad
        wavelength window. Cuts on per-role $A_{ki}$, the minimum
        branching ratio of the watched fluorescence channel
        (denominator = all decay channels of the final level), a
        same/different-wavelength detection constraint, a probe
        orbital filter, and optional isobaric-contaminant checks
        (an issue is flagged when a scheme laser can excite the
        isobar \emph{and} the isobar fluoresces near the detection
        wavelength). Results are scored
        ($\log_{10}$ weight $+ \sum \log_{10} A_{ki} +
        \log_{10}\mathrm{BR} - 5$ per contamination issue), deduped
        by physical pathway, and ranked; selecting a row draws the
        excitation/detection diagram and a full text description,
        and \app{Export results CSV\dots} writes the ranked list.
\end{description}

The browser's state --- spectrum, filters, search configuration, and
the ranked schemes (self-contained, so they restore offline) ---
persists in the \denis{} save file under the \code{nist\_asd} key.

\begin{figure}[hbtp]
    \centering
    \includegraphics[width=0.95\textwidth]{fig_tools_quick_plot.png}
    \caption{Quick Plot showing an ion-source mass-scan calibration
        loaded from CSV: ion current as a function of mass setting,
        with the sample isotopes appearing as well-resolved peaks
        near $A = 28$ and $A = 32$.}
    \label{fig:tools_quick_plot}
\end{figure}
